%%==========================================================================
\section{The Basic Tail Assignment Model}
\label{sec:basic}
%%==========================================================================

We restate the compact model of \citet{khaled2018compact} using the notation
of \Cref{tab:notation}.  This section covers the model without maintenance
constraints; they are introduced in \Cref{sec:maintenance}.

\begin{table}[t]
\centering\small
\caption{Notation summary.}
\label{tab:notation}
\begin{tabular}{ll}
\toprule
Symbol & Meaning \\
\midrule
$\F = \{1,\ldots,n_f\}$     & Set of flight legs \\
$\Pc = \{1,\ldots,n_p\}$    & Set of aircraft (tail numbers) \\
$\A = \{1,\ldots,n_k\}$     & Set of airports \\
$\aD{i}$, $\aI{i}$          & Departure / arrival airport of flight $i$ \\
$\tD{i}$, $\tI{i}$          & Departure / arrival time of flight $i$ (minutes) \\
$d_i$, $\bar d_i$           & Departure / arrival day derived from flight times \\
$\dt$                        & Minimum ground turn time (minutes) \\
$c_{ij}$                     & Assignment cost of flight $i$ to aircraft $j$ \\
$a^0_j$                      & Initial position airport of aircraft $j$ \\
$\Fl{k}{\theta}$             & $\{i \in \F : \aI{i} = k,\; \tI{i} \le \theta - \dt\}$ \\
$\Fd{k}{\theta}$             & $\{i \in \F : \aD{i} = k,\; \tD{i} < \theta\}$ \\
\bottomrule
\end{tabular}
\end{table}

\subsection{Decision variables and objective}

Let us define the decision variables and constraints for the basic tail assignment model.
\paragraph{C1 — Binary assignment domain.}
The decision variable $x_{ij}$ indicates whether flight $i$ is assigned to aircraft $j$.
\begin{equation}\label{eq:c0}
  x_{ij} = \begin{cases} 1 & \text{if flight } i \text{ is operated by aircraft } j, \\ 0 & \text{otherwise,} \end{cases}
  \qquad i \in \F,\; j \in \Pc.
\end{equation}

The objective minimises total assignment cost:
\begin{equation}\label{eq:obj}
  \min \sum_{i \in \F} \sum_{j \in \Pc} c_{ij} x_{ij}.
\end{equation}

\subsection{Constraints}

\paragraph{C2 — Flight coverage.}
Each flight must be assigned to exactly one aircraft:
\begin{equation}\label{eq:c1}
  \sum_{j \in \Pc} x_{ij} = 1, \qquad \forall i \in \F.
\end{equation}

\paragraph{C3,C4 — Equipment continuity (home/non-home airports).}
For any aircraft $j$, any non-home airport $k \ne a^0_j$, and any flight
$i$ departing from $k$, aircraft $j$ can be assigned to flight $i$ only if
it has already landed at $k$ early enough,
but at the home airport ($k = a^0_j$), one free unit is available at time zero
\begin{equation}\label{eq:c2}
  \sum_{i' \in \Fl{k}{\tD{i}}} x_{i'j}
  - \sum_{i' \in \Fd{k}{\tD{i}}} x_{i'j}
  \;\geq\; x_{ij},
  \qquad \forall j \in \Pc,\; k \in \A \setminus \{a^0_j\},\;
  i \in \F : \aD{i}=k.
\end{equation}

\begin{equation}\label{eq:c3}
  \sum_{i' \in \Fl{k}{\tD{i}}} x_{i'j}
  - \sum_{i' \in \Fd{k}{\tD{i}}} x_{i'j}
  \;\geq\; x_{ij}-1,
  \qquad \forall j \in \Pc,\; k = a^0_j,\;
  i \in \F : \aD{i}=k.
\end{equation}

\subsection{Validity}

Constraints~\eqref{eq:c2}--\eqref{eq:c3} are intended to encode both equipment
continuity (C2 of \citet{khaled2018compact}) and turn-time feasibility
(constraint~(6) of that paper).
The set $\Fl{k}{\tD{i}}$ counts only arrivals at least $\dt$ minutes before
the departure of flight $i$, so a prior arrival can supply flight $i$ only
after the required turn time. This mechanism is valid when departures are
strictly ordered, but it does not prevent two equal-time departures from using
the same available aircraft, as shown below.

\begin{proposition}[Claimed in \citealt{khaled2018compact}]
  Any feasible solution to \eqref{eq:c0}--\eqref{eq:c3}
  corresponds to a
  feasible flight plan satisfying the original continuity and turn-time
  conditions.
\end{proposition}

The proposition is not valid as stated when distinct flights may have equal
departure times. The following counterexample identifies the missing case.

\subsection{Counter-example with simultaneous departures}

The balance constraints~\eqref{eq:c2}--\eqref{eq:c3} use
$\Fd{k}{\theta}=\{i\in\F:\aD{i}=k,\;\tD{i}<\theta\}$ with a strict
inequality in departure time. If two flights leave the same airport at the
same time, neither flight appears in the other's departure-balance term.
This can make an assignment feasible while being physically impossible as a
flight plan.

\begin{example}[Feasible to \eqref{eq:c0}--\eqref{eq:c3}, infeasible in practice]
  \label{ex:simultaneous_departure_counterexample}
  Consider one aircraft $j$ with initial airport $a^0_j=\text{B}$ and
  turn time $\dt=30$ minutes. Let the three flights be:
  \begin{align*}
    &i_0: \text{B}\to\text{A},\quad \tD{i_0}=500,\;\tI{i_0}=550,\\
    &i_1: \text{A}\to\text{B},\quad \tD{i_1}=600,\;\tI{i_1}=700,\\
    &i_2: \text{A}\to\text{C},\quad \tD{i_2}=600,\;\tI{i_2}=700.
  \end{align*}
  Set $x_{i_0j}=x_{i_1j}=x_{i_2j}=1$.

  Coverage~\eqref{eq:c1} is satisfied (single aircraft, all flights assigned).
  Flight $i_0$ satisfies the home-airport constraint~\eqref{eq:c3}, because
  no flight has arrived at or departed from B before time 500. For $i_1$ at
  the non-home airport A:
  \begin{equation*}
    \Fl{\text{A}}{\tD{i_1}}=\{i_0\},\qquad
    \Fd{\text{A}}{\tD{i_1}}=\varnothing,
  \end{equation*}
  because $i_0$ departs from B and $\tD{i_2}=600\not<600$. Hence
  constraint~\eqref{eq:c2} gives
  \begin{equation*}
    \sum_{i'\in\Fl{\text{A}}{\tD{i_1}}}x_{i'j}
    -\sum_{i'\in\Fd{\text{A}}{\tD{i_1}}}x_{i'j}
    =1-0=1
    \ge x_{i_1j}=1.
  \end{equation*}
  Symmetrically for $i_2$, $\tD{i_1}=600\not<600$, so the same calculation gives
  $1\ge x_{i_2j}=1$. Therefore~\eqref{eq:c0}--\eqref{eq:c3} are feasible.

  However, the resulting plan is physically infeasible: one aircraft cannot
  operate both $i_1$ and $i_2$ simultaneously at time 600.
\end{example}

\paragraph{Implication.}
In implementations, this corner case is eliminated by adding pairwise
non-overlap constraints for time-overlapping flights on the same aircraft,
$x_{i_1j}+x_{i_2j}\le 1$.

\paragraph{Corrected C1--C5 feasibility criterion.}
Define $\mathcal{O}$ as the set of flight pairs that cannot be operated by one
aircraft because their operating intervals, including the turn-time buffer,
overlap:
\paragraph{Conflict-set definition.}
\begin{equation}
  \mathcal{O}=\left\{\{i_1,i_2\}\subseteq\F:
  [\tD{i_1},\tI{i_1}+\dt)\cap
  [\tD{i_2},\tI{i_2}+\dt)\neq\varnothing\right\}.
  \label{eq:overlap_set}
\end{equation}
For every pair in this conflict set, impose
\paragraph{C5 — Pairwise non-overlap.}
\begin{equation}\label{eq:c4_overlap}
  x_{i_1 j} + x_{i_2 j} \le 1,
  \qquad \forall \{i_1,i_2\}\in\mathcal{O},\; \forall j\in\Pc.
\end{equation}

The corrected core model is therefore
\eqref{eq:c0}--\eqref{eq:c4_overlap}.
This preserves the original $n_fn_p$ binary-variable space while restoring
physical executability. It does not, in general, preserve the linear
constraint count: the overlap family contains $|\mathcal{O}|n_p$ inequalities.

\subsection{Model size}

\begin{proposition}[Analogue of \citealt{khaled2018compact}, Proposition~2]
  The original model \eqref{eq:obj}--\eqref{eq:c3} with binary $x$ contains $n_f n_p$ binary
  decision variables and $O(n_fn_p)$ constraints, while the corrected model has $n_fn_p$ binary variables and $O(n_fn_p+|\mathcal{O}|n_p)$ constraints.
\end{proposition}

The corrected model retains $n_fn_p$ binary variables and has
$O(n_fn_p+|\mathcal{O}|n_p)$ constraints. In the worst case
$|\mathcal{O}|=O(n_f^2)$, although timetable sparsity normally makes the
conflict set much smaller. The original variable count compares favourably
with Levin's $O(n_f^2n_p)$-variable model; \citet{khaled2018compact} also
report that the time-space network formulation uses approximately
$1.25$--$2$ times as many variables and constraints on their test instances.

\subsection{Corrected proposition and its proof.}

The issue with the proof of feasibility in the original model is that it does not account for simultaneous departures 
of overlapping flights. Moreover, the logic of the published proof goes in the wrong direction, 
and instead of proving that any solution of C1--C4 defines correct 
feasibility, it proves that any feasible assignment flights satisfy C1-C4. 


The corrected proposition and its accompanying proof look like the following:

\begin{proposition}[Proposition~1]
  Any feasible solution to \eqref{eq:c0}--\eqref{eq:c4_overlap}
  corresponds to a feasible flight plan satisfying the original continuity and turn-time
  conditions.
\end{proposition}

\begin{proof}

  Let's formally prove it by induction on the number of planes.
  For the base case of one plane, any feasible assignment trivially satisfies the continuity and turn-time conditions:
\begin{itemize}
\item The single plane can start only for the airport it is assigned form the start, which is follows form the condition \ref{eq:c3}.
Let's define this assigned flight as $f_1$, so the time of departure and arrival are $t_{\text{dep}}(f_1)$ and $t_{\text{arr}}(f_1)$, respectively.
\item Any next $k$-th flight of this plane with time of departure $t_{\text{dep}}(f_k)$ must satisfy the turn-time condition $t_{\text{dep}}(f_k) \ge t_{\text{arr}}(f_{k-1}) + \tau$, where $\tau$ is the required turn time.
This follows directly from the turn-time condition \eqref{eq:c2} and overlap constraints \eqref{eq:c4_overlap}.
\end{itemize}
  This completes the base case.

  For the induction step, assume that any feasible assignment for $n-1$ planes satisfies the continuity and turn-time conditions. 
  Consider a feasible assignment for $n$ planes. 
  Select one plane and consider its sequence of assigned flights. 
  By the same argument as in the base case, this plane's flights satisfy the continuity and turn-time conditions. 
  Removing this plane from the assignment leaves a feasible assignment for the remaining $n-1$ planes, which by the induction hypothesis satisfies the continuity and turn-time conditions. Therefore, the entire assignment for $n$ planes satisfies the continuity and turn-time conditions.
  This completes the induction step.
\end{proof}






